\documentclass[11pt]{article}
\usepackage[margin=1in]{geometry}
\usepackage{amsmath,amssymb,amsthm,graphicx,booktabs}
\usepackage[hidelinks]{hyperref}
\newtheorem{theorem}{Theorem}
\newtheorem{remark}{Remark}
\newcommand{\D}{\mathbb D}
\newcommand{\C}{\mathbb C}
\newcommand{\BH}{\mathcal B(\mathcal H)}
\newcommand{\diag}{\operatorname{diag}}
\setlength{\parindent}{0pt}
\setlength{\parskip}{0.62em}

\title{No Universal Bohr Radius for Arbitrary\\Operator-Valued Biholomorphic Curves}
\author{Run \texttt{fa\_banach\_001}, lane 9}
\date{27 August 2026}

\begin{document}
\maketitle

\textbf{Status.} Candidate scoped counterexample, likely valid, for human
review.

\section{Source question}

Bhowmik and Das consider an operator-valued holomorphic function
\(f:\D\to\BH\) and a subordinate function
\[
 g(z)=\sum_{k=0}^{\infty}B_kz^k\in S(f).
\]
Their final Remark~(ii) asks whether there is a universal \(r_0>0\) such that
for every biholomorphic \(f\), every \(g\in S(f)\), and every
\(0\leq r\leq r_0\),
\begin{equation}\label{eq:bohr}
 \sum_{k=1}^{\infty}\|B_k\|r^k
 \leq \liminf_{|z|\to1^-}\|f(z)-f(0)\|.
\end{equation}
They suggest that the answer might be negative, perhaps even when \(f\) is
convex \cite[final Remark (ii), PDF p.~13]{BD}.

\begin{center}
\includegraphics[width=0.98\textwidth]{figures/open_problem_crop.png}
\end{center}

\section{Convention needed to read the question}

The paper's printed definition says that \(f(\D)\) is a domain in the ambient
space \(\BH\). If ``domain'' has its standard meaning of an ambient open set,
this is incompatible with every non-scalar example: the derivative of a map
from \(\C\) has rank at most one, so it cannot have a holomorphic local inverse
onto an open subset of a space of dimension greater than one. Thus, read
literally, the operator-valued part of the final question is empty.

The paper operates with the natural curve-level convention instead: \(f\) is
holomorphic and injective, and its inverse on the one-dimensional image is
holomorphic. Two internal signals make this intent concrete. The final
question treats \(A_1=f'(0)\) as an invertible member of \(\BH\), not as an
isomorphism \(\C\to\BH\). More decisively, the downloadable arXiv TeX contains
a commented-out explanatory paragraph that calls
\(z\mapsto\diag(z,t(z))\) biholomorphic for arbitrary holomorphic \(t\), in
order to observe that its higher coefficients can be arbitrarily large.

The following theorem answers the question under that operative convention.
It does not claim to satisfy the incompatible ambient-open clause.

\section{Claimed contribution}

\begin{theorem}\label{thm:main}
Under the source paper's operative curve convention, there is no universal
\(r_0>0\) for which \eqref{eq:bohr} holds for every operator-valued
biholomorphic \(f\) and every \(g\in S(f)\). A counterexample exists in
\(\mathcal B(\C^2)\) for every proposed positive radius.
\end{theorem}

\section{Proof intuition}

Keep one diagonal coordinate equal to \(z\). This permanently records the
parameter, so injectivity and a holomorphic inverse are automatic. Use the
other coordinate to add two arbitrarily large, cancelling Taylor
coefficients. Their norms add on the left of \eqref{eq:bohr}, but the
cancellation makes that coordinate return to the value one at the boundary.
The first coordinate simultaneously forces the boundary liminf to be exactly
one. Thus the coefficient side can be made larger than the geometric side at
any prescribed positive radius.

\section{Proof}

Let \(r_0>0\) be arbitrary and put
\[
 r=\min\{r_0/2,1/2\}.
\]
Then \(0<r\leq r_0\). Choose a real number
\begin{equation}\label{eq:Mchoice}
 M>\frac{1-r}{r^2(1+r)}
\end{equation}
and take \(\mathcal H=\C^2\). Define
\begin{equation}\label{eq:fdef}
 f_M(z)=
 \begin{pmatrix}
 z&0\\
 0&z+Mz^2-Mz^3
 \end{pmatrix},
 \qquad z\in\D,
\end{equation}
and let \(g=f_M\).

The map \(f_M\) is a holomorphic polynomial. Its first diagonal coordinate is
\(z\), so it is injective. Its inverse on \(f_M(\D)\) is the restriction of
the bounded complex-linear coordinate functional
\[
 \pi_{11}:\mathcal B(\C^2)\longrightarrow\C,
 \qquad \pi_{11}(T)=T_{11}.
\]
Consequently \(f_M\) is biholomorphic under the operative convention.
Taking the identity Schwarz function \(\phi(z)=z\) gives
\(g=f_M\circ\phi\), hence \(g\in S(f_M)\).

Since \(f_M(0)=0\), the nonzero Taylor coefficients of \(g\) are
\[
 B_1=I_2,\qquad
 B_2=\diag(0,M),\qquad
 B_3=\diag(0,-M).
\]
Their operator norms are \(1,M,M\), respectively. Therefore the left-hand
side of \eqref{eq:bohr} at the chosen radius is
\begin{align*}
 \sum_{k=1}^{\infty}\|B_k\|r^k
 &=r+Mr^2+Mr^3\\
 &=r+Mr^2(1+r)>1,
\end{align*}
where the last inequality is exactly \eqref{eq:Mchoice}.

It remains to compute the right-hand side. For every \(z\in\D\),
\[
 \|f_M(z)\|
 =\max\{|z|,\,|z+Mz^2-Mz^3|\}\geq |z|.
\]
It follows that
\[
 \liminf_{|z|\to1^-}\|f_M(z)\|\geq1.
\]
On the other hand, along positive real numbers \(z=\rho\uparrow1\),
\[
 \rho+M\rho^2-M\rho^3
 =\rho+M\rho^2(1-\rho)\longrightarrow1,
\]
while the first diagonal coordinate \(\rho\) also tends to one. Hence
\[
 \lim_{\rho\uparrow1}\|f_M(\rho)\|=1,
\]
and therefore
\begin{equation}\label{eq:rhs}
 \liminf_{|z|\to1^-}\|f_M(z)-f_M(0)\|=1.
\end{equation}
The strict coefficient estimate and \eqref{eq:rhs} contradict
\eqref{eq:bohr} at \(r\leq r_0\). Since \(r_0>0\) was arbitrary, no universal
positive radius exists. \qed

\section{Verification and limitations}

The script \texttt{code/verify\_counterexample.py} checks the coefficient
identities, the choice of \(M\), and the strict inequality with exact rational
arithmetic for representative radii. The proof itself is symbolic and uses no
numerical approximation.

The construction settles the unrestricted question only under the operative
curve convention. It is not convex. Accordingly it does not settle the
stronger suggestion that failure may persist for convex biholomorphic maps.
Under the printed ambient-open definition, the non-scalar class is empty
instead; that is a defect in the formulation, not an additional conclusion of
the counterexample.

\section{Bounded novelty check}

The search on 27 August 2026 covered the run's registry, solution, attempt,
and proof-gap indexes; exact searches for arXiv:1904.00351, its title, and
phrases from the final question; the arXiv source and Cambridge publisher
page; and the later operator-valued Bohr papers surfaced by those searches.
No explicit proof or counterexample to this unrestricted question was found.
This supports, but cannot certify, novelty.

\textbf{Human review recommendation.} Ask an expert in holomorphic mappings
in Banach spaces or operator-valued Bohr phenomena to confirm the source
convention. Once the operative convention is accepted, the mathematical
calculation is elementary and exact.

\begin{thebibliography}{9}
\bibitem{BD}
B. Bhowmik and N. Das,
\emph{Bohr phenomenon for operator valued functions},
Proc. Edinburgh Math. Soc. \textbf{64} (2021), 72--86;
arXiv:1904.00351,
\url{https://arxiv.org/abs/1904.00351};
\url{https://doi.org/10.1017/S0013091520000395}.
\end{thebibliography}

\end{document}
